\caption{\textbf{A detector-class visibility split does not transfer to matched TopK recovery.}
\textbf{a}, Pre-specified E4-v2 endpoint at $\lambda_{\mathrm{eff}}=0.55$: feature-independent detection power for axis and dense loadings. Thin lines pair the same seed across geometries; points are the eight seed-level observations, diamonds are means, and error bars are seed-level percentile-bootstrap $95\%$ confidence intervals (1,000 resamples). Power is $0.081$ $[0.073,0.090]$ for the axis loading and $1.000$ $[1.000,1.000]$ for the dense loading.
\textbf{b}, Decoder cosine from the same matched runs. Both geometries satisfy the frozen cosine-$0.80$ recovery criterion in all eight seeds (means $0.970$ and $0.978$ for axis and dense, respectively).
\textbf{c}, Detector--decoder decoupling over all 24 matched seed--strength pairs. Each point reports the dense-minus-axis difference in independence power and decoder cosine; pale trajectories connect the three strengths within a seed. The black star marks the pre-specified $\lambda_{\mathrm{eff}}=0.55$ mean ($\Delta_{\mathrm{ind}}=0.919$, $\Delta_{\mathrm{cos}}=0.008$).
\textbf{d}, Descriptive paired audit over the same 24 pairs. Detector class flips in $24/24$ pairs, whereas binary recovery status agrees in $22/24$ pairs. Panel d is descriptive rather than an additional pre-specified endpoint. The experiment uses eight seeds, $N=4096$ detector samples per trial, 400 alternative trials, a 16-atom TopK decoder with $k=2$, 20,000 training samples, and 4,000 optimization steps.}
\label{fig:detector_decoder_decoupling}
